\subsection{The ZX-calculus}
%%%% QUICK OVERVIEW
The ZX calculus~\cite{vdw_working_cs_zx} is a formal graphical language
composed of nodes and wires connecting them. Diagramatically, they are
represented as vertices on an undirected graph, alongside a \emph{phase}
in the interval $[0,2\pi)$. Using only these simple building blocks,
it is possible to construct a semantic meaning to encode linear maps
between quantum states. 

In its simplest form, every diagram can be written with two kinds of nodes,
usually called \emph{spiders}. They are differentiated between \emph{Z-spiders}
and \emph{X-spiders}, and are represented by green and red colored nodes
respectively. These spiders are interpreted as the following linear maps:

%%%%%%[Interpretación de las arañas zx]


%%%%%%%%% Hadamard generator
It is possible to incorporate to the calculus a third generator to represent the
Hadamard gate. This generator is not strictly necessary, since the operation can
be expressed in terms of red and green spiders. However, its inclusion to the
language greatly aides the ease of reading for diagrams. The generator is usually
drawn as a yellow box and is defined as follows:

%%%%%%%%%%%%% Definición de generador Hadamard


%%%%%%%%%%% [Cables identidad, cups y caps]

A ZX-diagram is an open graph, with wires acting as inputs on the left and outputs on
the right. The wires act as edges connecting the spiders, which are the nodes of the
graph. The spiders can be composed either sequentially ($\circ$) or in parallel
($\otimes$) which correspond to matrix and tensor products respectively. This yields
the global interpretation for the diagram. Take for example:

%%%%%%% Ejemplo simple de diagrama

The core idea behind these diagrams is that \emph{only the connectivity matters}, and
it is possible to freely deform the wires or move the nodes without changing the meaning.

The main strength of working with ZX diagrams is incorporating an abstraction layer over
the mathematics needed to compute large and intricate linear maps from smaller blocks.
It leverages this layer with a set of rewriting rules which allow to simplify the diagrams,
while keeping the same semantics overall. The rules are the following:

%%%%%%%%%%%%%% ZX rules
 
While several different equalities can be derived, the previously presented rules form a
\emph{complete} set. This means, that if two different ZX-diagrams represent the same linear
map, then there exists a sequence of rule applications such that one diagram is transformed
into the other.

\iffalse
A \zxGND\ diagram is generated by the following primitives,
in addition to parallel and serial composition:
\[
    \tikzfig{szx/elem/spiderZ} : n_1 \to m_1
    \qquad
    \tikzfig{szx/elem/spiderX} : n_1 \to m_1
\]
\[
    \tikzfig{szx/elem/hadamard} : 1_1 \to 1_1
    \qquad
    \tikzfig{szx/elem/ground} : 1_1 \to 0_1
\]
\[
    \tikzfig{szx/elem/wire} : 1_1 \to 1_1
    \qquad
    \tikzfig{szx/elem/cup} : 0_1 \to 2_1
\]
\[
    \tikzfig{szx/elem/cap} : 2_1 \to 0_1
    \qquad
    \tikzfig{szx/elem/swap} : 2_1 \to 2_1
    \qquad
    \tikzfig{szx/elem/empty} : 0_1 \to 0_1
\]
\fi